\chapter{迭代器}
\label{ch:17}

迭代器把“接下来还有什么”封装成状态，从而让遍历、判断和归约可以组合。它和原容器不是同一个对象：容器中仍有数据，不代表已经消费过的迭代器还能从头读取。

本章先用 all 与 any 显示消费位置，再用 fold 追踪累加器，最后手动处理 Set 与 Map 的 next。理解状态之后，再使用更长的组合链，才不容易把重复调用和重新遍历混淆。

\section{消费状态与提前结束}

先看一个容易误读的场景：同一个列表先检查“是否全为正数”，再检查“是否存在大于 2 的数”。如果两次都从列表开头遍历，两项都应成立；但复用同一个迭代器时，第二次从上次停下的位置继续。

all 与 any 都在消费这个位置状态，而不是只保存一条待执行的条件。下面把复用 first 与重新创建 second 并排比较，让第二次结果为何不同可以直接从执行过程推出来。

\par
\medskip
\begin{lstlisting}[style=chali]
import std::iter;

void main()
{
    List<int64> xs;
    xs.push(1);
    xs.push(2);
    xs.push(3);
    xs.push(4);
    let first = xs.iter();
    println("第一次 all: {}", first.all(fn(int64 const& x)
    {
        x > 0
    }));
    println("第二次 any: {}", first.any(fn(int64 const& x)
    {
        x > 2
    }));
    let second = xs.iter();
    println("新建后 any: {}", second.any(fn(int64 const& x)
    {
        x > 2
    }));
}
\end{lstlisting}
\par
\medskip

\par
\medskip
\begin{lstlisting}[style=out]
第一次 all: 1
第二次 any: 0
新建后 any: 1
\end{lstlisting}
\par
\medskip

\subsection*{代码详解}

first 是有当前位置的迭代器。all 对每项判断大于 0，列表中 1 到 4 都满足，所以返回真，同时把该迭代器消费到末尾。紧接着对 first 调 any 时已没有剩余项，因此返回假，即使原列表中本来存在大于 2 的元素。

second 来自重新调用 xs.iter，拥有新的遍历位置，any 于是能找到 3 或 4，返回真。消费的是迭代状态，不是删除 xs 的元素；也不要由此认为 all 一定遍历全部输入，遇到不满足条件时它可以提前结束。本例全为正数才到达末尾。

\section{从序列建立管道}

先明确输入序列，再组合变换，最后通过收集或汇总取得结果。不要把构造管道与执行完整遍历混为一谈；闭包捕获的有效期必须覆盖迭代器实际使用的时间。

\par
\medskip
\begin{lstlisting}[style=chali]
import std::iter;

void main()
{
    List<int64> xs;
    xs.push(2);
    xs.push(4);
    xs.push(6);
    int64 total = fromList<int64>(xs).fold(0, fn(int64 acc, int64 x)
    {
        acc + x
    });
    println("{} {}", total, xs.len());
}
\end{lstlisting}
\par
\medskip

\par
\medskip
\begin{lstlisting}[style=out]
12 3
\end{lstlisting}
\par
\medskip

\subsection*{代码详解}

fromList<int64>(xs) 为归约提供元素序列，fold 的第一个参数 0 是初始累加值。闭包的 acc 接收上轮累计结果，x 接收当前元素，每次返回 acc+x，依次得到 2、6、12。

fold 的结果交给 total，原列表仍有三项，所以输出 \texttt{12 3}。累加器类型与闭包返回类型需要一致；若把加法改成别的运算，还应重新选择合适的初值，不能认为任何 fold 都从 0 开始。

\section{Map 和 Set 的遍历}

Map 的迭代项包含键和值，Set 的迭代项是键。业务若需要固定顺序，应显式排序或选择有序结构；不要把一次运行的哈希表顺序写成持久化格式。

\par
\medskip
\begin{lstlisting}[style=chali]
import std::iter;

void main()
{
    Set<int64> s;
    s.add(1);
    s.add(2);
    s.add(3);
    let si = s.iter();
    int64 sum = 0;
    loop
    {
        match(si.next())
        {
            Value(v) => { sum = sum + v; }
            None => { break; }
        }
    }
    println("sum = {} len = {}", sum, s.len());
    Map<int64, int64> m;
    m.put(1, 10);
    m.put(2, 20);
    let mi = m.iter();
    int64 doubled = 0;
    loop
    {
        match(mi.next())
        {
            Value(pair) => { doubled = doubled + pair.first * pair.second; }
            None => { break; }
        }
    }
    println("doubled = {} len = {}", doubled, m.len());
}
\end{lstlisting}
\par
\medskip

\par
\medskip
\begin{lstlisting}[style=out]
sum = 6 len = 3
doubled = 50 len = 2
\end{lstlisting}
\par
\medskip

\subsection*{代码详解}

第一段手动调用集合迭代器的 next，每次用 Option 表示“还有一项”或“已经结束”。Value(v) 时累加，None 时 break，最后集合元素之和为 6，原集合长度仍为 3。这个求和不依赖遍历顺序。

Map 的 next 返回键值对，pair.first 为键，pair.second 为值。本例累计的是 1*10 加 2*20，即 50，并非只把值相加，变量名 doubled 也不改变实际公式。两段都只读容器；遍历期间插入、删除等结构修改需要遵守迭代器约束，不能因为长度在本例不变就认为任意修改都安全。

\section{结构修改与失效}

遍历期间的扩容、删除或重排可能使当前位置和元素借用失效。遵守迭代器提供的修改方式；需要批量修改时，可以先收集稳定标识再在第二阶段更新。提前退出循环仍要清理迭代器和临时拥有值。

\subsection*{小结}

一次遍历可分为取得起点、逐项推进、处理元素和结束。all、any、fold 会消费当前状态，next 用 Option 显式呈现是否还有元素；需要重来时从容器取得新的迭代器。

遍历是否改变原容器、元素是值还是借用、何种修改使迭代失效，都要按具体接口确认。不要把某个顺序无关的求和结果当成稳定遍历顺序的证明。下一章把程序组织与模块依赖纳入视野。

\section{习题}

\begin{enumerate}

\item 连续两次调用 any 时会重新从头开始吗？

\item 何时真正执行一个迭代管道？

\item Map 遍历顺序能否直接当作排序？

\item 如何安全地批量删除匹配元素？

\end{enumerate}

\subsection*{参考解答与验收要点}

\begin{enumerate}

\item 不能假定；检查消费状态或重建迭代器。

\item 由消费或终止操作推进，按具体协议核对。

\item 不能，需要时显式排序。

\item 用受支持接口，或先收集标识再修改，避免保留失效引用。

\end{enumerate}

下一章：\hyperref[ch:18]{第 18 章}。

\section{番外：少了一页的遍历}

改写自《西游记》通天河经书落水、晒经的情节。

却说老鼋将师徒并经担掀入水中，师徒上岸，拣高处晒经。三藏看那经卷湿透，心疼不已。八戒自告奋勇：“老猪最会遍历，晒一本收一本，绝不漏项。”

他指着第一本，记数，收走；经书向前一挪，他的索引也向前一加。如此行了半日，收得倒快，只是每隔一本便有一本仍在地上滴水。

三藏问其故，八戒道：“弟子每次都递增，难道上进也错？”悟空扯过书担：“你向前一步，队伍也向前一步，一来一去，正好跨过一位师兄。”

沙僧于是先遍历记下要收的书，再另作一轮搬运。不可结构修改的迭代器，便老老实实守它的约定。八戒嫌分两轮不好看，又在纸上写了一条十分漂亮的惰性链，写罢便去吃斋。

到天黑，经书仍湿。悟空道：“你那张纸还没有消费者呢。”

三藏揭起《本行经》，不慎粘破了尾页，叹曰：“此乃此道不全。”八戒忙道：“不关遍历的事，这是缓存粘连。”悟空笑道：“先把你跳过的那几本捡回来，再给故障换名字不迟。”

\subsection*{技术对照}

惰性迭代器通常在消费时执行工作。修改当前元素与改变集合结构不是同一操作；遍历中的引用、源集合存活和结构修改按具体接口规则处理。
